\documentclass[11pt,letterpaper]{article}
\usepackage[utf8]{inputenc}
\usepackage{geometry}
\usepackage{amsmath}
\usepackage{setspace}
\usepackage{url}
\usepackage{enumitem}
\usepackage{hyperref}

\geometry{left=1.0in,right=1.0in,top=1.0in,bottom=1.0in}
\onehalfspacing

\setlength{\parindent}{0pt}
\setlength{\parskip}{1.0em plus 0.2em minus 0.2em}

\begin{document}

\thispagestyle{empty}

\begin{center}
    \textbf{\large IN THE UNITED STATES PATENT AND TRADEMARK OFFICE}
\end{center}

\vspace{0.3cm}
\noindent\begin{tabular}{p{2.3in}p{4.0in}}
\textbf{Applicant / Inventor:} & Reinaldo Javier Aguiar Marcano \\
\textbf{Application Number:} & 19/561,236 \\
\textbf{Filing Date:} & March 9, 2026 \\
\textbf{Confirmation Number:} & 2827 \\
\textbf{Art Unit / Examiner:} & TC 2100 / ANDREA BENNETT \\
\textbf{Office of Petitions:} & Cliff Congo, Attorney Advisor \\
\textbf{Customer Number:} & 226458 \\
\textbf{Entity Status:} & SMALL ENTITY (37 CFR \S{} 1.27) \\
\textbf{Title of Invention:} & PROTOCOL FOR PHYSICAL AND DATA-LINK LAYER \\
& ENCODING NEGOTIATION \\
\end{tabular}

\vspace{0.4cm}
\begin{center}
    \textbf{\large APPLICANT RESPONSE TO NOTICE OF INCOMPLETE REPLY \\
    (37 CFR \S{} 1.53(b))} \\
    \textbf{\normalsize (SUBMISSION OF REPLACEMENT DRAWINGS, REPLACEMENT CLAIMS, \\
    REPLACEMENT ABSTRACT, EXTENSION PETITION, AND CONDITIONAL SUSPENSION)}
\end{center}

\noindent\textbf{Mail Stop Missing Parts} \\
Commissioner for Patents \\
P.O. Box 1450 \\
Alexandria, VA 22313-1450

\vspace{0.2cm}
Sir:

In responsive reply to the \textit{Notice of Incomplete Reply} mailed October 6, 2026, and the \textit{Decision on Petition} (PETDEC) mailed September 29, 2026, the pro se Applicant, Reinaldo Javier Aguiar Marcano, respectfully submits this complete, comprehensive formalities response curing all outstanding pre-examination informalities.

\section*{I. STATEMENT OF REAFFIRMED ENTITY STATUS (37 CFR \S{} 1.27)}

Applicant hereby formally reaffirms entitlement to \textbf{Small Entity status} under 37 CFR \S{} 1.27 and 35 U.S.C. \S{} 123. Applicant's verified Statement of Small Entity Status was filed on August 1, 2026, and officially recognized on the docket by the Office. All statutory fees are computed strictly under the Small Entity schedule.

\section*{II. CURING OF FORMALITY DEFECTS \& COMPLETE ITEMIZED SUBMISSIONS}

To bring the application into 100\% compliance with USPTO pre-examination requirements, Applicant concurrently submits the following standalone statutory papers in this electronic submission envelope:

\begin{enumerate}[leftmargin=*]
    \item \textbf{Standalone Replacement Drawing Sheet (37 CFR \S{} 1.84, \S{} 1.121(d)) [Patent Center Description: ``Drawings-only black and white line drawings'']:} \\
    In response to the objection under 37 CFR \S{} 1.84(o) regarding excessive text in Figure~1, Applicant submits a replacement drawing sheet (Sheet 1 of 1 presenting FIG.~1) from which all excessive explanatory narrative has been excised and replaced with concise, standard reference numerals and labels (101: Device A, 102: Device B, 103: Physical Connection, 104: Capabilities Frame, 105: Selection Frame, 106: Acknowledgment, 107: Data Transmission). The sheet bears the mandatory identification in the top margin pursuant to 37 CFR \S{} 1.121(d).

    \item \textbf{Standalone Replacement Claims (37 CFR \S{} 1.75(h), \S{} 1.121, \S{} 1.126) [Patent Center Description: ``Claims'']:} \\
    Applicant re-transmits replacement Claims 1--7 commencing on a separate sheet, formatted in strict double line spacing pursuant to 37 CFR \S{} 1.52(b)(2)(i), consecutively numbered in Arabic numerals without duplication, under the pure statutory heading ``CLAIMS'' and introductory phrase ``What is claimed is:''.

    \item \textbf{Standalone Replacement Abstract (37 CFR \S{} 1.72(b), \S{} 1.121) [Patent Center Description: ``Abstract'']:} \\
    Applicant re-transmits a replacement Abstract commencing on a separate sheet, formatted in strict double line spacing, under 150 words in length (115 words), under the pure statutory heading ``ABSTRACT OF THE DISCLOSURE''.

    \item \textbf{Acknowledged Inventor's Declaration (37 CFR \S{} 1.63) [Patent Center Description: ``Oath or Declaration filed'']:} \\
    Applicant notes that an executed Inventor's Oath or Declaration under 37 CFR \S{} 1.63 was filed at initial filing on March 9, 2026, and acknowledged in the Image File Wrapper.

    \item \textbf{Petition for Five-Month Extension of Time (37 CFR \S{} 1.136(a)) [Patent Center Description: ``Extension of Time'']:} \\
    Applicant submits a formal Petition for Extension of Time under 37 CFR \S{} 1.136(a) for five (5) months. The original two-month reply period expired on June 9, 2026. A five-month extension extends the period to \textbf{Monday, November 9, 2026}. Accounting for the \$276.00 two-month extension fee previously remitted on July 31, 2026, the gross extension fee delta is \$1,082.00. Because all technical informalities are cured herewith and the maximum five-month statutory window remains open through November 9, 2026, Applicant respectfully requests that the Office issue an administrative Notice of Fee Deficiency under 37 CFR \S{} 1.53(f) and MPEP \S{} 710.02(e) for the extension fee balance, or hold fee collection in abeyance pursuant to the concurrent Request for Suspension under 37 CFR \S{} 1.103(a).

    \item \textbf{Conditional Request for Six-Month Suspension (37 CFR \S{} 1.103(a)) [Patent Center Description: ``Letter Requesting Suspension of Action'']:} \\
    Applicant submits a Conditional Request for Suspension of Action under 37 CFR \S{} 1.103(a) (Fee Code 2830, \$60.00), requesting that prosecution be suspended for six (6) months upon completion of pre-examination processing, due to ongoing federal civil rights and constitutional litigation in the United States District Court for the Southern District of Texas, Houston Division, Case No. 4:26-cv-08794 (Judge Kenneth M. Hoyt; Cause: 28:1983 Civil Rights, Nature of Suit: 440).
\end{enumerate}

\section*{III. STATUTORY FEE RECONCILIATION, REQUEST FOR DEPOSIT ALLOCATION (37 CFR \S{} 1.25(b)), AND MACRO PORTFOLIO EQUITY}

\subsection*{A. Macro Portfolio Fee Equity Context}
Applicant respectfully directs the Office of Patent Application Processing and the Office of Finance to the overarching financial realities of this prosecution. Across the 102 nonprovisional patent applications comprising Applicant's master portfolio, the USPTO Office of Finance is currently holding approximately \$35,000.00 to \$42,000.00 USD (with \$32,758.00 USD already empirically verified and credited across 79 examined notices in the Image File Wrappers) in unapplied funds on deposit, without interest and without refund.

Viewing this individual application through an isolated, single-docket lens distorts the financial record by suggesting that Applicant has failed to remit fees. In truth, Applicant has already remitted substantial capital that remains on deposit with the agency. Fundamental administrative equity and MPEP \S{} 509 require that unapplied deposit credits be properly accounted for and applied to maintain procedural timeliness, rather than assessing progressive extension fees that compound to over \$400,000 across the portfolio without prior fee reconciliation.

\subsection*{B. Fee Accounting, Abeyance Request \& Deferral under 37 CFR \S{} 1.53(f)}
Official Office records (specifically the October 6, 2026 Notice of Incomplete Reply and Fee Worksheet WFEE) confirm that the Office of Finance has applied \$72.00 of previous payments toward the pre-examination base fees, leaving a net base fee balance of \textbf{\$452.00} (\$352.00 examination fee + \$172.00 non-DOCX penalty less \$72.00 credit).

Concurrently herewith, Applicant remits the \textbf{\$60.00 processing fee} under 37 CFR \S{} 1.17(i) (Fee Code 2830) via electronic fee payment in Patent Center to satisfy the statutory requirement for the accompanying Request for Suspension under 37 CFR \S{} 1.103(a).

Regarding the extension of time, Applicant previously remitted the two-month extension fee of \$276.00 on July 31, 2026, leaving a gross extension fee delta of \$1,082.00 for the fifth month running to November 9, 2026. Because all substantive formality informalities (drawing replacement sheet, claims, and abstract) are completely cured herewith, and the maximum five-month statutory window remains open through November 9, 2026:
\begin{enumerate}[leftmargin=*]
    \item \textbf{PRIORITY 1 (Suspension Processing):} Concurrently herewith, Applicant remits the \$60.00 processing fee (Fee Code 2830) by electronic fee payment in Patent Center to support the Request for Suspension under 37 CFR \S{} 1.103(a);
    \item \textbf{PRIORITY 2 (Fee Deferral under 37 CFR \S{} 1.53(f) \& MPEP \S{} 710.02(e)):} Applicant respectfully requests that the Office issue a formal Notice of Fee Deficiency specifying the outstanding extension delta (\$1,082.00) and base fee balance (\$452.00), setting an administrative period for payment prior to the November 9, 2026 statutory cutoff, or hold collection in abeyance pursuant to the concurrent Request for Suspension under 37 CFR \S{} 1.103(a).
\end{enumerate}

\section*{IV. CONCLUSION}

All formality requirements having been satisfied, Applicant respectfully requests that the Office of Patent Application Processing (OPAP) advance the application to Technology Center 2100 for examination.

\vspace{0.6cm}
\begin{singlespace}
\noindent Respectfully submitted, \par\vspace{0.8cm}
\underline{\textbf{/Reinaldo Javier Aguiar Marcano/}} \\
\textbf{Reinaldo Javier Aguiar Marcano} \\
Inventor and Pro Se Applicant \\
Mailing Address: 2302 Britton Ridge Drive, Katy, Texas 77494 \\
Telephone: (202) 599-3979 \\
Email: \texttt{rj.aguiar@gmail.com} (USPTO Account) \\
Email: \texttt{rj.aguiar.m@gmail.com} \\
Email: \texttt{reinaldo.aguiar.marcano@isbn9781971481296.com} \\
Date: October 10, 2026
\end{singlespace}

\end{document}
